% sections/theory.tex -- method and theory (SPEC.md sections 1--4).
% Uses macros.tex. Results: Theorem 1 (exactness), Corollary 1.1 (rebuild),
% Lemma 2 (apparent pairs), Corollaries 2.1--2.2 (Morse order), Proposition 3
% (cost), Proposition 4 (no worst-case locality). Executable checks of every
% statement: research/mmhm/reference/check_props.py.

\section{Method and theory}\label{sec:theory}

This section defines the data structure maintained by \method{mmhm}, proves
that it yields exact Betti numbers after every atomic edit
(Theorem~\ref{thm:exact}), shows where discrete Morse theory enters and what it
provably buys (Lemma~\ref{lem:apparent} and its corollaries), states the cost of
an update in terms of quantities of the data structure
(Proposition~\ref{prop:cost}), and shows that these quantities admit no bound in
terms of the size of the update (Proposition~\ref{prop:nonlocal}). Correctness
does not depend on any locality or on the order used; the order affects only
cost. Coefficients are in $\Z$ throughout; we make no statement about integer
homology or torsion. The arguments hold in every dimension; the implementation
is restricted to $\dim K\le 3$.

\subsection{Setting}\label{sec:setting}

Let $K$ be a finite simplicial complex of dimension $d$. A $k$-simplex is a set of
$k+1$ vertices; $\Kc{k}$ is the set of $k$-simplices of $K$ and $n_k=|\Kc{k}|$.
We identify the chain group $\Cc{k}=C_k(K;\Z)$ with the set of finite subsets of
$\Kc{k}$, addition being symmetric difference. The boundary
$\bd_k\colon\Cc{k}\to\Cc{k-1}$ sends a simplex to the set of its facets
(codimension-one faces); $\Bm{k}$ is its matrix, with rows indexed by
$\Kc{k-1}$ and columns by $\Kc{k}$. We write $\Zc{k}=\ker\bd_k$,
$r_k=\rank\Bm{k}$ for $1\le k\le d$, and $r_0=r_{d+1}=0$. By rank--nullity,
$\dim\Zc{k}=n_k-r_k$ and $\dim\operatorname{im}\bd_{k+1}=r_{k+1}$, so
\begin{equation}\label{eq:betti}
  \beta_k \;=\; n_k - r_k - r_{k+1}\qquad(0\le k\le d).
\end{equation}

\paragraph{Edits and updates.} An \emph{atomic edit} is either
$\Ins(\sigma)$, for a simplex $\sigma\notin K$ all of whose facets lie in $K$, or
$\Del(\sigma)$, for a simplex $\sigma\in K$ that has no proper coface in $K$.
Each atomic edit leaves a simplicial complex. An \emph{update} $u_t$ is a finite
sequence of atomic edits (a batch); its \emph{size} is $s_t=|u_t|$. The task is
to output $\beta_0,\dots,\beta_d$ of $K_t$ after every update $t$.

\begin{observation}[Structural fact]\label{obs:structural}
Let $\sigma$ be a $k$-simplex. While $\sigma\in K$, its boundary $\bd\sigma$ does
not change. Inserting $\sigma$ appends the column $\bd\sigma$ to $\Bm{k}$ and
appends a zero row (indexed by $\sigma$) to $\Bm{k+1}$; deleting $\sigma$
removes its column from $\Bm{k}$ and removes its row from $\Bm{k+1}$, and that
row is zero. No other boundary matrix changes, and among $n_0,\dots,n_d$ only
$n_k$ changes.
\end{observation}

\begin{proof}
The boundary of a simplex is determined by its vertex set. When $\sigma$ is
inserted it has no coface: a coface would already be in $K$, and $K$ is closed
under faces while $\sigma\notin K$. A deleted simplex has no coface by the
precondition of $\Del$. In both cases the row $\sigma$ of $\Bm{k+1}$ is zero.
\end{proof}

Consequently, for each $k\ge1$ the problem is to maintain $r_k=\rank\Bm{k}$ under
\emph{column insertion and column deletion}, where rows appear and disappear
only while they are zero; adding or removing a zero row changes no rank. The
vertex dimension $k=0$ needs no matrix: $\bd_0=0$.

\subsection{The dynamic reduced basis}\label{sec:drb}

Every cell $x\in K$ carries an \emph{order key} $\ord(x)$; keys of distinct cells
are distinct. For a non-zero chain $c$, $\low(c)$ is the cell of $c$ with the
largest key. Keys do not change between rebuilds (Section~\ref{sec:morse}); an
inserted cell receives a key larger than every key in use. Nothing in this
subsection depends on how the keys are chosen.

\begin{definition}[DRB]\label{def:drb}
Let $k\ge1$. A \emph{dynamic reduced basis} of $\Bm{k}$ is a finite family of
\emph{slots} $j\in J_k$, each holding a pair $(R_j,V_j)$ with $V_j\in\Cc{k}$
and $R_j\in\Cc{k-1}$, such that
\begin{description}
\item[(I1)] $R_j=\bd_k V_j$ for every $j\in J_k$;
\item[(I2)] the non-zero $R_j$ have pairwise distinct lows;
\item[(I3)] $(V_j)_{j\in J_k}$ is a basis of $\Cc{k}$; in particular $|J_k|=n_k$.
\end{description}
Two index structures are derived from the slots: $\pivot(\ell)=j$ if $R_j\ne0$
and $\low(R_j)=\ell$ (well defined by (I2)), and, for $x\in\Kc{k}$,
$\occ(x)=\{j\in J_k : x\in V_j\}$, the \emph{footprint} of $x$. Both are updated
together with the chains; we do not spell out this bookkeeping.
\end{definition}

Writing $R$ and $V$ for the matrices with columns $R_j$ and $V_j$, (I1) and (I3)
say that $R=\Bm{k}V$ with $V$ invertible, and (I2) says that $R$ is
\emph{reduced}. Unlike the decomposition used for persistence, $V$ is not
required to be triangular with respect to $\ord$.

\paragraph{Insertion.} $\Ins(\sigma)$ for a $k$-simplex $\sigma\notin K$, $k\ge1$,
all of whose facets are in $K$:
\begin{enumerate}
\item give $\sigma$ a key larger than all keys in use; create a slot $j$ with
  $V_j=\{\sigma\}$ and $R_j=\bd\sigma$;
\item while $R_j\ne0$ and $p=\pivot(\low(R_j))$ is defined: $R_j\gets R_j+R_p$,
  $V_j\gets V_j+V_p$;
\item if $R_j\ne0$, set $\pivot(\low(R_j))=j$.
\end{enumerate}
The slots $p_1,\dots,p_m$ added in step~2 form the \emph{reduction chain} of
$\sigma$. We call $\sigma$ \emph{positive} if $R_j=0$ at the end and
\emph{negative} otherwise.

\paragraph{Deletion.} $\Del(\sigma)$ for $\sigma\in\Kc{k}$, $k\ge1$, without
cofaces. Let $O=\occ(\sigma)$ and $Z=\{j\in O: R_j=0\}$.
\begin{description}
\item[Case A ($Z\ne\emptyset$).] Choose any $z\in Z$ (the implementation takes
  one with $|V_z|$ minimal). For every $j\in O\setminus\{z\}$:
  $V_j\gets V_j+V_z$. Remove slot $z$.
\item[Case B ($Z=\emptyset$).] Let $p\in O$ be the slot whose $\low(R_p)$ has
  the smallest key. For every $j\in O\setminus\{p\}$: $R_j\gets R_j+R_p$ and
  $V_j\gets V_j+V_p$. Remove slot $p$ and the entry $\pivot(\low(R_p))$.
\end{description}
Finally remove $\sigma$ from $K$. Inserting or deleting a vertex changes only
$n_0$.

\begin{theorem}[Exactness]\label{thm:exact}
Let $k\ge1$ and let $(R_j,V_j)_{j\in J_k}$ satisfy (I1)--(I3) for $\Bm{k}(K)$,
for some injective key function.
\begin{enumerate}
\item[(a)] $r_k=\#\{j\in J_k : R_j\ne0\}$, and $\{V_j : R_j=0\}$ is a basis of
  $\Zc{k}(K)$.
\item[(b)] Let $\sigma\notin K$ be a $k$-simplex whose facets are all in $K$.
  The loop of $\Ins(\sigma)$
  terminates, and afterwards the slots satisfy (I1)--(I3) for
  $\Bm{k}(K\cup\{\sigma\})$; every family satisfying (I1)--(I3) for
  $\Bm{k+1}(K)$ satisfies them, unchanged, for $\Bm{k+1}(K\cup\{\sigma\})$.
\item[(c)] Let $\sigma\in\Kc{k}$ have no coface. Then $O\ne\emptyset$; after
  $\Del(\sigma)$, in Case A and in Case B, the slots satisfy (I1)--(I3) for
  $\Bm{k}(K\setminus\{\sigma\})$; every family satisfying (I1)--(I3) for
  $\Bm{k+1}(K)$ satisfies them, unchanged, for $\Bm{k+1}(K\setminus\{\sigma\})$.
\item[(d)] Inserting a positive $k$-simplex increases $\beta_k$ by one; inserting
  a negative one decreases $\beta_{k-1}$ by one. A deletion in Case A decreases
  $\beta_k$ by one; a deletion in Case B increases $\beta_{k-1}$ by one. No
  other Betti number changes. Inserting (deleting) a vertex increases
  (decreases) $\beta_0$ by one. Moreover, $\Del(\sigma)$ is in Case A if and only
  if $\sigma$ lies in the support of a $k$-cycle of $K$, and $\Ins(\sigma)$ is
  positive if and only if $\sigma$ lies in the support of a $k$-cycle of
  $K\cup\{\sigma\}$.
\end{enumerate}
\end{theorem}

\begin{proof}
(a) By (I1) and (I3), $R=\Bm{k}V$ with $V$ invertible, so $\rank R=r_k$. The
non-zero columns of $R$ are linearly independent: in a non-empty sum of them,
let $R_i$ be the summand whose low $\ell$ has the largest key; every other
summand has its low, hence all its cells, strictly below $\ell$, so $\ell$
survives in the sum. Hence $r_k$ is the number of non-zero $R_j$. Each $V_j$ with
$R_j=0$ is a cycle; these chains are independent as part of a basis, and there
are $n_k-r_k=\dim\Zc{k}$ of them.

(b) Let $K'=K\cup\{\sigma\}$. Throughout the loop, $R_j=\bd V_j$ (true at the
start and preserved by linearity because $R_p=\bd V_p$) and $V_j=\sigma+w$ with
$w\in\Cc{k}(K)$. An addition is performed only when $\low(R_p)=\low(R_j)=\ell$;
all cells of both chains have keys at most $\ord(\ell)$ and $\ell$ cancels, so
afterwards $R_j=0$ or $\low(R_j)$ is strictly smaller. Hence the loop
terminates, after at most one addition per non-zero slot. On exit either
$R_j=0$ or $\low(R_j)$ is no pivot, so recording the pivot preserves (I2); (I1)
holds. Since $\Cc{k}(K')=\Cc{k}(K)\oplus\Z\sigma$ and the old $V_i$ form a basis
of $\Cc{k}(K)$, the old chains together with $V_j=\sigma+w$ form a basis of
$\Cc{k}(K')$: they are $n_k+1$ chains spanning $\sigma=V_j+w$. For $\Bm{k+1}$:
the chain group $\Cc{k+1}$ is unchanged, each $R_i=\bd V_i$ is supported on facets
of $(k+1)$-simplices of $K'$, none of which is $\sigma$
(Observation~\ref{obs:structural}), and keys of existing cells are unchanged;
so (I1)--(I3) continue to hold.

(c) Since the $V_j$ span $\Cc{k}$, $\sigma$ is a sum of some of them, so some
$V_j$ contains $\sigma$: $O\ne\emptyset$. Each operation $V_j\gets V_j+V_q$ with
$j\ne q$ is an invertible change of basis, so throughout the loop the family
remains a basis of $\Cc{k}(K)$. For $j\in O\setminus\{q\}$ (with $q=z$ in Case A,
$q=p$ in Case B), both $V_j$ and $V_q$ contain $\sigma$, so $V_j+V_q$ does not;
slots outside $O$ never contained $\sigma$. After the loop, $V_q$ is therefore
the only chain containing $\sigma$. Removing slot $q$ leaves $n_k-1$ independent
chains, all in $\Cc{k}(K\setminus\{\sigma\})$, which has dimension $n_k-1$:
(I3) holds.
\emph{Case A.} Because $\bd V_z=R_z=0$, the additions leave every $R_j$
unchanged, so (I1) and (I2) hold; the removed slot had $R_z=0$ and owned no
pivot.
\emph{Case B.} Every $R_j$ with $j\in O$ is non-zero, and by (I2) the lows are
distinct, so $p$ is well defined and $\low(R_p)$ is strictly smaller than
$\low(R_j)$ for $j\in O\setminus\{p\}$. All cells of $R_p$ have keys at most
$\ord(\low(R_p))<\ord(\low(R_j))$, so $\low(R_j)\notin R_p$ and
$\low(R_j+R_p)=\low(R_j)$; in particular $R_j+R_p\ne0$. Lows are unchanged and
the low of $R_p$ is released with slot $p$, so (I2) holds; (I1) holds by
linearity. (The minimal choice of $p$ is what guarantees this: adding an $R_p$
with a larger low would move lows and require further reduction.)
For $\Bm{k+1}$: $\sigma$ has no coface, so no $R_i=\bd V_i$ of dimension $k+1$
contains $\sigma$ and $\sigma$ is not a pivot there; deleting the zero row
$\sigma$ leaves (I1)--(I3) intact.

(d) By (a), $r_k$ increases by one exactly when the inserted $\sigma$ is negative,
and decreases by one exactly in Case B; $n_k$ changes by $\pm1$; $r_{k+1}$ is
unchanged because only a zero row of $\Bm{k+1}$ is added or removed
(Observation~\ref{obs:structural}); all other $n_i$ and $r_i$ are unchanged.
Formula~\eqref{eq:betti} gives the stated changes, also for vertices. If
$z\in Z$, then $V_z$ is a cycle containing $\sigma$. Conversely, if a cycle
$c\in\Zc{k}(K)$ contains $\sigma$, write $c$ in the basis
$\{V_j:R_j=0\}$ of $\Zc{k}(K)$ from (a); some $V_j$ in this expansion contains
$\sigma$, and that $j$ lies in $Z$. For insertion, if $\sigma$ is positive then
the final $V_j$ is a cycle of $K'$ containing $\sigma$. Conversely, if a cycle
$c$ of $K'$ contains $\sigma$, then $\bd\sigma=\bd(c+\sigma)$ lies in the
column space of $\Bm{k}(K)$, so $r_k$ does not increase, and by (a) the new slot
has $R_j=0$.
\end{proof}

A batch is processed one atomic edit at a time, so Theorem~\ref{thm:exact}
makes the Betti numbers exact after every update. The choice of $z$ in Case A
is free; it only affects sparsity. Theorem~\ref{thm:exact} holds for every key
function, so the order chosen in Section~\ref{sec:morse} is a matter of cost
alone.

\paragraph{Initialisation and rebuilds.}
The DRB is initialised, and periodically rebuilt, by the static reduction with
clearing~\cite{chen2011twist}, tracking $V$: for $k=d,d-1,\dots,1$, process the
columns of $\Bm{k}$ in increasing key order; if the $k$-simplex $\tau$ is the low
of a non-zero reduced column $R$ of $\Bm{k+1}$, \emph{clear} it ($R_\tau:=0$,
$V_\tau:=R$); otherwise reduce it as in steps~2--3 of $\Ins$, starting from
$V_\tau=\{\tau\}$, $R_\tau=\bd\tau$.

\begin{corollary}[Rebuild]\label{cor:rebuild}
For every injective key function, the static reduction with clearing produces
slots satisfying (I1)--(I3) in every dimension. Moreover, every $V_\tau$ consists
of $\tau$ and cells with smaller keys. In particular the ranks it reports are
exact even when $V$ is not stored.
\end{corollary}

\begin{proof}
We show $\low(V_\tau)=\tau$ for every column $\tau$, by induction along the
processing order. A non-cleared $\tau$ ends with
$V_\tau=\{\tau\}+\sum V_p$ over pivot owners $p$ processed earlier, hence with
smaller keys, and by induction every cell of $V_p$ has key at most $\ord(p)$. A
cleared $\tau$ has $V_\tau=R$ with $\low(R)=\tau$. Ordering rows and columns by
key, $V$ is therefore upper unitriangular, hence invertible: (I3). (I1) holds
for non-cleared columns by linearity, and for cleared ones because $R=\bd V'$
for the slot $V'$ of $\Bm{k+1}$ that produced it (induction over dimensions), so
$\bd V_\tau=\bd\bd V'=0=R_\tau$. (I2): a non-cleared column ends with $R_\tau=0$ or
with a low that no earlier column owns; cleared columns are zero. The last
statement follows from Theorem~\ref{thm:exact}(a).
\end{proof}

Note that the correctness of clearing here follows from the invariants alone;
we do not need the classical fact that a cleared column would have reduced to
zero. After deletions, however, $V$ is no longer triangular (Case A adds
$V_z$, whose low may exceed that of $V_j$), so the DRB maintains ranks and a cycle
basis but not a persistence pairing, and we claim none. This is the difference
from vineyard updates~\cite{cohensteiner2006vines}, which keep $V$ triangular
for a changing order.

\paragraph{Rebuild policy.} \method{mmhm} rebuilds with the Morse order of Section~\ref{sec:morse}; it
rebuilds after update $t$ if $t-t_{\mathrm{last}}\ge m$ or
$\nnz(R)+\nnz(V)>(1+\alpha)\max(\nnz_{\mathrm{last}},|K|)$, where
$\nnz_{\mathrm{last}}$ is the value right after the previous rebuild (defaults
$m=\infty$, $\alpha=1$), and the rebuild is charged to update $t$.
\method{drb-id} uses insertion keys and never rebuilds; \method{mmhm-norb} uses
the Morse order at the initial build only.

\subsection{The Morse order}\label{sec:morse}

Discrete Morse theory~\cite{forman2002user} enters only through the choice of
keys at a rebuild. A rebuild computes an acyclic matching by collapses, in the
spirit of~\cite{benedetti2014random}, and orders cells by reverse removal time.

\begin{definition}[Collapse matching and Morse order]\label{def:collapse}
Run the following on $P:=K$, with a counter $t:=0$ and a queue $Q$ of the cells
$x$ that have exactly one immediate coface in $P$ (\emph{free faces}).
\begin{enumerate}
\item While $Q\ne\emptyset$: pop $x$; skip it unless $x\in P$ and $x$ has
  exactly one immediate coface $y$ in $P$. Then \emph{pair} $(x,y)$, set
  $\rmt(y):=t$, $\rmt(x):=t+1$, $t\gets t+2$, remove $y$ and $x$ from $P$, and
  push every facet of $y$ or $x$ that has become free.
\item If $P\ne\emptyset$, choose a cell $c$ of the highest dimension present in
  $P$ (smallest vertex tuple), declare it \emph{critical}, set $\rmt(c):=t$,
  $t\gets t+1$, remove it, push the facets that have become free, and go to~1.
\end{enumerate}
The \emph{Morse order} is $\ord(x):=-\rmt(x)$ (relabelled to $0,\dots,|K|-1$).
Write $M$ for the set of pairs, $\Crit$ for the critical cells, $\Lo$ for the
cells $x$ with $(x,y)\in M$ and $\Up$ for the cells $y$ with $(x,y)\in M$; for
$x\in\Lo$ let $V(x)$ be its partner. For a cell $x$, its \emph{youngest facet}
is the facet with the largest key and its \emph{oldest coface} the coface with
the smallest key.
\end{definition}

The removed cells are always maximal in $P$, and $P$ remains a subcomplex. For
the paired cell $y$ this uses the simplicial diamond property: if $y$ had a
coface $w\in P$, then $w$ would have exactly two facets containing $x$, namely
$y$ and some $y'\ne y$, and $y'\in P$ because $P$ is a subcomplex, so $x$ would
not be free. The queue discipline is not specified; all results below hold for
every discipline.

A \emph{gradient path} (V-path) is a sequence $\tau_0,\sigma_0,\tau_1,\sigma_1,
\dots,\tau_m$ with $(\tau_i,\sigma_i)\in M$ and $\tau_{i+1}$ a facet of
$\sigma_i$ different from $\tau_i$.

\begin{lemma}[Gradient pairs are apparent pairs]\label{lem:apparent}
For every execution of the collapse procedure, with $\ord=-\rmt$:
\begin{enumerate}
\item[(a)] $\ord$ is a filtration order: $\ord(f)<\ord(s)$ whenever $f$ is a facet
  of $s$.
\item[(b)] $\rmt$ strictly increases along the lower cells of every gradient
  path ($\rmt(\tau_{i+1})>\rmt(\tau_i)$); in particular $M$ has no closed
  gradient path, so it is an acyclic matching.
\item[(c)] For every $(\tau,\sigma)\in M$: $\tau$ is the youngest facet of
  $\sigma$, i.e.\ $\tau=\low(\bd\sigma)$, and $\sigma$ is the oldest coface of
  $\tau$.
\item[(d)] If $z\ne0$ is a cycle, then $\low(z)\notin\Up$.
\item[(e)] In the static reduction under $\ord$, with or without clearing, for
  every $(\tau,\sigma)\in M$ the column $\sigma$ is not cleared, receives no
  addition, and obtains pivot $\tau$; thus $R_\sigma=\bd\sigma$ and
  $(\tau,\sigma)$ is an \emph{apparent pair}~\cite{bauer2021ripser}.
\end{enumerate}
\end{lemma}

\begin{proof}
(a) Only maximal cells of $P$ are removed (above), so every cell is removed
before each of its facets.

(b) When $\sigma_i$ is removed at time $\rmt(\sigma_i)$, its facet $\tau_{i+1}$ is
still in $P$. Times $\rmt(\sigma_i)$ and $\rmt(\tau_i)=\rmt(\sigma_i)+1$ are
used by $\sigma_i$ and $\tau_i$, and $\tau_{i+1}\ne\tau_i$, so
$\rmt(\tau_{i+1})>\rmt(\tau_i)$. A closed gradient path would give
$\rmt(\tau_0)<\rmt(\tau_0)$.

(c) When $\sigma$ is removed, all its facets are in $P$; $\sigma$ and $\tau$ take
the times $t$ and $t+1$, and the other facets are removed later, so they have
smaller keys than $\tau$. When $\tau$ was popped, $\sigma$ was its only coface in
$P$; every other coface of $\tau$ had been removed earlier, so it has a larger
key than $\sigma$.

(d) Let $u=\low(z)\in\Up$ with $u=V(w)$. The coefficient of $w$ in $\bd z=0$
vanishes and $w\in\bd u$, so some $u'\in z$ with $u'\ne u$ has $w$ as a facet.
By (c), $u$ is the oldest coface of $w$, so $\ord(u')>\ord(u)$, contradicting
$u=\low(z)$.

(e) If $\sigma$ were cleared, it would be the low of a reduced column $R$ of the
next dimension; $R=\bd V'$ (Corollary~\ref{cor:rebuild}) is a non-zero cycle with
low $\sigma\in\Up$, contradicting (d). Columns are processed in increasing key
order. When $\sigma$ is processed, every earlier column $x$ has $R_x=0$ (if
cleared) or $R_x=\bd V_x$ with all cells of $V_x$ of key at most
$\ord(x)<\ord(\sigma)$ (Corollary~\ref{cor:rebuild}); by (c) no such cell has
$\tau$ as a facet, so $\tau\notin R_x$ and $\pivot(\tau)$ is undefined. Since
$\low(\bd\sigma)=\tau$ by (c), the reduction loop performs no addition and
assigns $\pivot(\tau)=\sigma$.
\end{proof}

\begin{corollary}[Work happens only on critical columns]\label{cor:critical-work}
In the static reduction with clearing under the Morse order, the column of every
$x\in\Lo$ with $\dim x\ge1$ is cleared, the column of every $y\in\Up$ is
apparent, and therefore every column addition is performed while reducing the
column of a critical cell.
\end{corollary}

\begin{proof}
For $(x,y)\in M$ with $\dim y=k+1\ge2$, Lemma~\ref{lem:apparent}(e) applied to
$\Bm{k+1}$, which is processed before $\Bm{k}$, gives $R_y=\bd y$ with
$\low(R_y)=x$, so $x$ is cleared. Upper cells are apparent by
Lemma~\ref{lem:apparent}(e). Every remaining column belongs to a critical cell.
\end{proof}

\paragraph{Morse boundary.} By Lemma~\ref{lem:apparent}(b) the following
recursion is well founded and defines a linear map
$\flow\colon\Cc{k-1}\to\Cc{k-1}$: $\flow(\tau)=\tau$ for $\tau\in\Crit$,
$\flow(\tau)=0$ for $\tau\in\Up$, and
$\flow(\tau)=\sum_{\tau'\in\bd V(\tau),\,\tau'\ne\tau}\flow(\tau')$ for
$\tau\in\Lo$. Thus $\flow(\tau)$ is the sum, over all gradient paths from $\tau$
to a critical cell, of their endpoints, counted mod~2. The \emph{Morse boundary}
of a critical $k$-cell $c$ is $\dM c:=\flow(\bd c)$, the boundary of the Morse
complex of the matching~\cite{forman2002user}; the Morse boundary matrices
$\Bm{k}^M$ have rows $\Crit\cap\Kc{k-1}$ and columns $\Crit\cap\Kc{k}$. For
$y\in\Up$ with partner $x$, the definition of $\flow(x)$ gives
\begin{equation}\label{eq:flow-kills}
  \flow(\bd y)=\flow(x)+\sum_{f\in\bd y,\,f\ne x}\flow(f)=0 .
\end{equation}

\begin{corollary}[Critical columns perform the Morse reduction]\label{cor:morse-reduction}
Consider the static reduction with clearing under the Morse order and, next to
it, the standard reduction with clearing of the Morse boundary matrices
$\Bm{d}^M,\dots,\Bm{1}^M$, critical cells taken in increasing key order. For
every critical $k$-cell $c$, $k\ge1$:
\begin{enumerate}
\item[(i)] the additions performed on column $c$ are first additions of apparent
  columns $R_y=\bd y$ ($y\in\Up$), then additions of columns of critical cells;
\item[(ii)] if $c$ is not cleared, then when the first phase ends the column
  contains no cell of $\Lo$ and its restriction to $\Crit$ equals $\dM c$;
\item[(iii)] the second phase performs, one for one, the additions that the
  Morse reduction performs on column $c$; $c$ is cleared exactly when it is
  cleared in the Morse reduction; and the final column $R_c$ contains no cell
  of $\Lo$, while $R_c\cap\Crit$ is the reduced Morse column of $c$.
\end{enumerate}
\end{corollary}

Each first-phase addition $R\gets R+\bd y$ with $\low(R)=x$ and $y=V(x)$
replaces $x$ by the other facets of $y$; it is one step of the discrete gradient
flow, and $\flow(R)$ does not change by~\eqref{eq:flow-kills}. This is the
precise sense in which the reduced column of a critical cell ``traces its
gradient flow''. $R_c$ may still contain cells of $\Up$; we observed this in
the randomised checks and claim nothing about them.

\begin{proof}
We use the rule of Definition~\ref{def:collapse} that a critical cell is chosen
of the highest dimension present in $P$. (Lemma~\ref{lem:apparent} did not use
it.)

\emph{Claim (separation).} For every $q$, every critical $q$-cell has a smaller
key than every $q$-cell of $\Lo$. Indeed, when a critical $q$-cell $t$ is removed,
$P$ contains no $(q+1)$-cell. If $\ell\in\Lo$ were a $q$-cell with
$\rmt(\ell)>\rmt(t)$, its partner $V(\ell)$, a $(q+1)$-cell removed at time
$\rmt(\ell)-1\ge\rmt(t)$, would be in $P$ at that moment.

We argue by induction over the dimensions (top-down) and, within $\Bm{k}$, over
the columns in increasing key order. The statement about clearing is shown below;
a cleared column receives no addition and has $R_c=0$, matching the zero column
of the Morse reduction. So let $c$ be a column that is not cleared, and let $R$
denote its column during its reduction; initially $R=\bd c$.

\emph{First phase.} Suppose $\ell=\low(R)\in\Lo$ and let $y=V(\ell)$. All cells
added to $V_c$ so far are $c$ and upper cells with smaller keys, so $\ell$ is a
facet of a cell $x$ with $\ord(x)\le\ord(c)$; since $y$ is the oldest coface of
$\ell$ (Lemma~\ref{lem:apparent}(c)) and $y\ne c$, $\ord(y)<\ord(c)$. By
Lemma~\ref{lem:apparent}(e), the column of $y$ has been processed, holds
$R_y=\bd y$ and owns $\pivot(\ell)$; the algorithm adds it, which
strictly decreases the low and preserves $\flow(R)$ by~\eqref{eq:flow-kills}.
The phase ends with $R=0$ or with $t=\low(R)\notin\Lo$. In the latter case $R$
is a non-zero cycle (by (I1), $R=\bd V_c$), so $t\notin\Up$ by
Lemma~\ref{lem:apparent}(d), and $t\in\Crit$. By the separation claim every
cell of $R$, having key at most $\ord(t)$, lies outside $\Lo$. As $\flow$ fixes
critical cells and kills upper cells,
$R\cap\Crit=\flow(R)=\flow(\bd c)=\dM c$. This proves (ii) and the first part
of (i).

\emph{Clearing.} The cells cleared in $\Bm{k}$ are the lows of the non-zero
reduced columns of $\Bm{k+1}$: the lower cells (lows of apparent columns) and
the lows of non-zero critical columns. By the induction hypothesis (iii) for
dimension $k+1$, the latter are critical and coincide with the lows of the
non-zero reduced Morse columns, i.e.\ with the cells cleared in the Morse
reduction. So a critical $c$ is cleared in one reduction if and only if in the
other; in the top dimension nothing is cleared in either.

\emph{Second phase.} A pivot on a critical row cannot be owned by an apparent
column (its low is in $\Lo$) or a cleared one (zero), so it is owned by a
critical column $c'$ processed earlier. By the induction hypothesis $R_{c'}$ has
no cell of $\Lo$, it is a cycle, its low is critical and equals
$\low(R_{c'}\cap\Crit)$, and $R_{c'}\cap\Crit$ is the reduced Morse column of
$c'$; hence the pivot tables of the two reductions agree on critical rows.
Adding such columns keeps $R\cap\Lo=\emptyset$ and keeps $R$ a cycle, so while
$R\ne0$ its low is critical and equals $\low(R\cap\Crit)$. The sequence of pivot
look-ups and additions on $R\cap\Crit$, which starts from $\dM c$, is therefore
exactly the Morse reduction of column $c$, and no apparent column is added any
more. This proves (i) and (iii).
\end{proof}

\begin{remark}[Relation to known results]\label{rem:morse-literature}
Bauer~\cite{bauer2021ripser} observed that the apparent pairs of any simplexwise
filtration form a discrete gradient and used them to skip columns.
Lemma~\ref{lem:apparent} is a converse for the gradient produced by collapses:
ordering cells by reverse removal time turns every gradient pair into an
apparent pair. Mischaikow and Nanda~\cite{mischaikow2013morse} showed that
Morse reductions compatible with a filtration preserve persistent homology;
Corollary~\ref{cor:morse-reduction} is a matrix-level form of this for our order:
on critical rows and columns, the static reduction \emph{is} the reduction of
the Morse complex, and by Corollary~\ref{cor:critical-work} the other columns
cost no additions. This is what the \method{full-morse} baseline exploits. In the
language of algebraic Morse theory, the first phase of
Corollary~\ref{cor:morse-reduction} is Gaussian elimination of the invertible
entries $\langle\bd y,x\rangle=1$ of the matched pairs, which is how the Morse
complex arises algebraically; the order makes the reduction algorithm perform
these eliminations by itself, without a separate Morse-reduction pass.
\end{remark}

\begin{remark}[Scope]\label{rem:morse-scope}
Lemma~\ref{lem:apparent} and its corollaries describe the state produced by a
rebuild. Between rebuilds the matching is not maintained, inserted cells get
keys larger than all others, and deletions add chains whose lows may exceed
those of their targets. We make no claim that the DRB stays sparse between
rebuilds; the Morse order is a heuristic for small footprints and short
reduction chains, its effect is measured, and the rebuild policy bounds the
drift of $\nnz(R)+\nnz(V)$.
\end{remark}

\subsection{Cost of an update}\label{sec:cost}

We count \emph{elementary operations}: inserting, removing or testing one cell in
a chain or one slot in a footprint, one pivot-table operation, and one low query.
With chains stored as ordered sets keyed by $\ord$ (or as hash sets with a
separately maintained maximum), each costs $O(\log n)$ time, $n=|K|$.

\begin{proposition}[Output-sensitive cost]\label{prop:cost}
Let $\sigma$ be a $k$-simplex, $k\ge1$.
\begin{enumerate}
\item[(a)] $\Ins(\sigma)$ with reduction chain $p_1,\dots,p_m$ performs
  $\Theta\bigl(k+1+\sum_{i=1}^{m}(1+|R_{p_i}|+|V_{p_i}|)\bigr)$ elementary
  operations. The $p_i$ are pairwise distinct, so $m\le r_k$ and
  $\sum_i(|R_{p_i}|+|V_{p_i}|)\le\nnz(R)+\nnz(V)$ of the DRB of $\Bm{k}$.
\item[(b)] $\Del(\sigma)$ with footprint $O=\occ(\sigma)$ performs
  $\Theta\bigl(|O|\,(1+|V_z|)\bigr)$ elementary operations in Case A and
  $\Theta\bigl(|O|\,(1+|R_p|+|V_p|)\bigr)$ in Case B, including the maintenance
  of $\occ$ and $\pivot$.
\item[(c)] An update costs the sum over its atomic edits; a vertex edit costs
  $O(1)$.
\end{enumerate}
\end{proposition}

\begin{proof}
(a) Each addition touches every cell of $R_p$ and $V_p$ once and performs one
pivot look-up and one low query; the final chain $V_j$, of size at most
$1+\sum_i|V_{p_i}|$, is entered into the footprints; $|\bd\sigma|=k+1$. Along
the chain the lows strictly decrease (proof of Theorem~\ref{thm:exact}(b)), and
each slot owns one low, so no slot is added twice.
(b) Scanning $O$ to determine $Z$, and $z$ or $p$, costs $O(|O|)$ (sizes $|V_j|$
are stored). In Case A each of the $|O|-1$ additions touches the $|V_z|$ cells of
$V_z$ in $V_j$ and in the footprints, and removing slot $z$ removes it from
$|V_z|$ footprints. Case B is the same with $|R_p|+|V_p|$, plus one pivot
removal.
\end{proof}

The quantities in Proposition~\ref{prop:cost} --- footprints $|\occ(\sigma)|$,
reduction-chain lengths and the sizes of the chains involved --- depend on the
whole history and on the order, not on $s_t$. We report them as measurements and
state no locality theorem. Proposition~\ref{prop:nonlocal} shows that none is
possible for this data structure.

\subsection{No worst-case locality}\label{sec:nonlocal}

\begin{proposition}[No worst-case locality for the DRB]\label{prop:nonlocal}
Let $K$ be either (W5a) the cycle graph $C_L$, $L\ge3$, and $\sigma$ one of its
edges, with $n:=L$; or (W5b) a triangulated $2$-sphere (more generally, a
connected closed triangulated surface) and $\sigma$ one of its triangles, with
$n:=n_2$. Let $k=\dim\sigma$. Then, for every key function and every choice in
Case A:
\begin{enumerate}
\item[(a)] for every DRB of $\Bm{k}(K)$ satisfying (I1)--(I3), $\Del(\sigma)$ is in
  Case A with $|V_z|=n$ and performs at least $n$ elementary operations;
\item[(b)] for every DRB of $\Bm{k}(K\setminus\{\sigma\})$ satisfying
  (I1)--(I3), $\Ins(\sigma)$ is positive, ends with $V_j=\Kc{k}$, and its
  reduction chain satisfies $\sum_i|V_{p_i}|\ge n-1$;
\item[(c)] if $S$ is any DRB state of $K\setminus\{\sigma\}$ and $S'$ is the state
  after $\Ins(\sigma)$, then $\Del(\sigma)$ applied to $S'$ restores $S$ exactly.
  Hence, after the first deletion, toggling $\sigma$ alternates between two fixed
  states $S_0,S_1$, and each toggle costs at least $n$ and at most
  $O\bigl(n+\nnz(R)+\nnz(V)\bigr)$ elementary operations, with $\nnz$ taken in
  $S_0$.
\end{enumerate}
So on these families every update of size $s_t=1$ costs this algorithm
$\Omega(n)$ work, and $\Theta(n)$ whenever $\nnz(R)+\nnz(V)=O(n)$ in $S_0$.
\end{proposition}

The statement bounds the work of \emph{this} algorithm, which must maintain
(I1)--(I3) and the footprint index; it is not a lower bound for computing Betti
numbers under edits. For (W5a), $\beta_1=n_1-n_0+\beta_0$ and $\beta_0$ can be
maintained by dynamic connectivity in polylogarithmic amortised time per edge
edit~\cite{holm2001polylog}; for (W5b), see Remark~\ref{rem:surfaces}. What the
proposition does show is that any method that stores cycle representatives
explicitly, as (I3) forces, has to write the $n$-cell cycle when it appears.

\begin{proof}
In both cases $\Zc{k}(K)=\{0,\Kc{k}\}$. For $C_L$, a $1$-cycle is a set of edges
meeting every vertex an even number of times, and the only non-empty such set is
the whole cycle. For a connected closed surface, a $2$-cycle that contains a
triangle contains, for each edge of it, the other triangle on that edge; as the
surface is strongly connected, it contains every triangle. Also
$\Zc{k}(K\setminus\{\sigma\})=\{0\}$, since the only non-zero cycle of $K$
contains $\sigma$.

(a) By Theorem~\ref{thm:exact}(a) exactly one slot $z$ has $R_z=0$, and
$V_z=\Kc{k}\ni\sigma$. So $z\in O$, the deletion is in Case A with $Z=\{z\}$, and
removing slot $z$ removes it from the footprints of all $n$ cells of $V_z$.

(b) $\beta_k(K)=\beta_k(K\setminus\{\sigma\})+1$, so by
Theorem~\ref{thm:exact}(d) $\sigma$ is positive; the final $V_j$ is then a non-zero
cycle of $K$, i.e.\ $V_j=\Kc{k}$. Since $V_j=\{\sigma\}+\sum_iV_{p_i}$, we get
$n\le1+\sum_i|V_{p_i}|$.

(c) In $S$ no chain contains $\sigma$, and $\Ins(\sigma)$ creates one slot $j$
and modifies no other. So in $S'$ we have $\occ(\sigma)=\{j\}$ and $R_j=0$ by
(b): $\Del(\sigma)$ is in Case A with $O=Z=\{j\}$, modifies nothing and removes
$j$. The key given to $\sigma$ differs between insertions, but it is always the
largest key in use and the algorithms only compare keys, so every insertion from
$S_0$ performs the same computation. The bounds follow from (a), (b) and
Proposition~\ref{prop:cost}.
\end{proof}

\begin{remark}[Surfaces need no linear algebra]\label{rem:surfaces}
Let $S$ be a connected triangulated surface, possibly with boundary (every edge
lies in one or two triangles and every vertex link is a path or a cycle). Over
$\Z$: $\beta_0=1$; $\beta_2=1$ if no edge lies in exactly one triangle and
$\beta_2=0$ otherwise (by strong connectivity a non-zero $2$-cycle is the set of
all triangles, whose boundary is the set of boundary edges); and
$\beta_1=\beta_0+\beta_2-\chi(S)$ with $\chi=n_0-n_1+n_2$. Under atomic edits that
keep $S$ a connected surface, $n_0,n_1,n_2$ and the number of boundary edges are
maintained in $O(1)$ per edit from per-edge triangle counts. All Betti numbers
are then available in $O(1)$ per edit without any matrix --- provided that
connectivity and the surface property are known in advance; the counters do not
certify connectivity. This is the \method{euler} method. Consequently, workloads
consisting of connected surfaces with known connectivity cannot evaluate the
algebraic part of a homology-maintenance method: a method that switches to these
formulas returns correct Betti numbers there whether or not its matrix updates
are correct. The gates of version~1 (its \S2.4) are these formulas, and its
experiments P1 and P3 are workloads of this type
(Section~\ref{sec:erratum}).
\end{remark}
